\section{实施路线与软件边界}

当前正式实现只保留一条主链：
\begin{equation}
 g
 \rightarrow A_{\rm em}(g,T),B(g)
 \rightarrow \{D^{-1},A_{\rm em}\}\text{ operator polynomial}
 \rightarrow \mathcal C_\theta(A,R)
 \rightarrow \mathrm{FGMRES}
 \rightarrow X
 \rightarrow q_{\rm Joule}
 \rightarrow a(t).
\end{equation}
其中 Maxwell 始终在固定背景完整 edge space 中求解；不存在全局 Maxwell solution basis、Maxwell rank、dense $Q/S$ 特征或 reduced-coefficient 网络输出。神经网络不预测完整 Maxwell correction，而只为真实 operator 生成的短阶 polynomial directions 预测少量复数混合系数。最终电磁结果必须满足真实 sparse residual 条件
\begin{equation}
 \frac{\lVert B-A_{\rm em}X\rVert_2}{\lVert B\rVert_2}\le \varepsilon_{\rm EM}.
\end{equation}

\subsection{核心模块}
正式统一路径的职责如下：
\begin{verbatim}
sdfmpneo/
  unified_geometry.py         # analytic / implicit geometry
  unified_background.py       # fixed multiscale edge/cell physics
  unified_neural_operator.py  # MR-Jacobi + operator polynomial coefficients
  unified_maxwell.py          # true-residual FGMRES closure
  unified_dataset.py          # geometry/material residual sampling metadata
  unified_trainer.py          # solution-label-free directional training + benchmark
  unified_thermal.py          # residual-driven automatic thermal basis
  unified_model.py            # exact Joule + thermal dynamics + save/load
  unified_runtime.py          # single train/predict runtime
\end{verbatim}
历史研究模块可以保留用于追溯，但不得由顶层 API、\texttt{run.py} 或正式 CLI 自动切换成另一条模型路线。

\subsection{Maxwell polynomial 预条件器}
对当前真实 residual $r$，先形成 Jacobi 方向 $q_0=D^{-1}r$，并求最优复数标量 $\alpha$：
\begin{equation}
 \alpha=\arg\min_{\beta\in\mathbb C}\lVert r-\beta A_{\rm em}q_0\rVert_2.
\end{equation}
因此零网络对应的物理 baseline
\begin{equation}
 z_{\rm MR}=\alpha D^{-1}r
\end{equation}
不会使一步 residual 比原 residual 更差。

随后只用真实 $A_{\rm em}$ 与 $D^{-1}$ 构造短阶 Krylov-like polynomial directions：
\begin{equation}
 q_{k+1}=D^{-1}A_{\rm em}q_k.
\end{equation}
各方向按 RMS 归一化。默认三阶预条件器写成
\begin{equation}
 \boxed{
 z_\theta
 =z_{\rm MR}
 +\sum_{k=0}^{2}c_k(\mathcal F(A_{\rm em},r);\theta)\widehat q_k
 }.
\end{equation}
网络只输出三个有界复数系数 $c_k$；空间方向完全来自当前真实物理算子。

\subsection{Maxwell 训练流程}
训练只学习如何组合真实 operator directions，以更快地降低 Maxwell residual：
\begin{enumerate}
  \item 在几何与物理材料温升分布中采样 $(g,T)$；
  \item 由固定背景装配真实 sparse $A_{\rm em}(g,T)$ 与端口 RHS $B(g)$；
  \item 从端口 RHS、端口线性组合以及基础物理迭代产生的中间 residual 构造训练 residual bank；
  \item 网络读取当前 residual、MR-Jacobi correction、MR 后 residual、operator diagonal/phase、sparse row coupling、edge orientation 与物理位置；
  \item 真实算子生成 $\widehat q_k$，网络只预测 $c_k$；
  \item 训练目标与 FGMRES 对预条件方向的实际使用方式一致：
  \begin{equation}
    L_{\rm EM}
    =\min_{\alpha\in\mathbb C}
    \frac{\lVert R-\alpha A_{\rm em}z_\theta\rVert_2^2}
         {\lVert R\rVert_2^2};
  \end{equation}
  \item 推理时把同一 polynomial smoother 作为 FGMRES 的可变预条件器，且只由真实 residual 决定停止。
\end{enumerate}

训练缓存不保存 Maxwell 解标签，也不保存 dense reduced matrices。网络最后一层零初始化，因此未训练网络严格退化到 MR-Jacobi；神经输出出现非有限值时，同一次 correction 也退化到 MR-Jacobi，而不会返回 surrogate Maxwell 答案。

\subsection{solver-level 验证}
Directional loss 只衡量单个预条件方向，不能替代真实 Krylov 性能。因此每次正式训练结束后，必须在 validation/test 子集上直接比较 MR-Jacobi baseline 与 neural polynomial preconditioner，并记录：
\begin{itemize}
  \item FGMRES iteration 的 median、90th percentile 与 maximum；
  \item restart 统计；
  \item final true residual 与成功率；
  \item wall time 与 wall-time speedup。
\end{itemize}
只有这些 solver-level 指标才能说明神经模块是否真的产生端到端加速。

\subsection{Joule 与 thermal ROM}
只有通过 Maxwell true-residual 检查后的完整场 $X$ 才允许进入 Joule 与阻抗计算。海水等体积导电区域使用
\begin{equation}
 q^{\rm volume}=\frac12\sigma |E|^2,
\end{equation}
线圈 AC resistance heating 作为物理项加入。

thermal ROM 保留，因为热扩散空间通常可压缩；thermal basis 的 rank 由真实 Joule anchor 与热方程 residual 自动决定，而不是手工固定。时间演化继续求解
\begin{equation}
 M_r(g)\dot a=-K_r(g)a+q_r(a,g,u),
\end{equation}
时间不是神经网络输入。

\subsection{禁止的耦合}
正式实现中不允许：
\begin{itemize}
  \item 用全局 Maxwell solution basis 或手工 Maxwell rank 作为神经载体；
  \item 用网络直接预测最终 Maxwell 场、阻抗、Joule heat 或温度；
  \item 把时间作为黑箱网络输入绕过 thermal ODE；
  \item 让网络学习 $M_r(g)$ 或 $K_r(g)$；
  \item 用 domain probe、Gate、fast/general/legacy 模式切换另一套模型；
  \item 在未满足真实 Maxwell residual 时静默返回近似答案。
\end{itemize}

\subsection{验证与正确性边界}
正式数值检查至少覆盖：
\begin{itemize}
  \item 零初始化 polynomial network 与 MR-Jacobi baseline 的一致性；
  \item polynomial directions 的尺度有限且来自真实 operator action；
  \item 非有限神经输出不会污染最终物理解；
  \item FGMRES 返回结果满足真实 full-space sparse residual；
  \item training report 含 baseline 与 neural 的 solver-level benchmark；
  \item Joule、阻抗只使用 corrected Maxwell field；
  \item thermal rank 由 residual 目标自动决定。
\end{itemize}

小 algebraic residual 只控制离散线性系统求解误差，并不自动保证连续 Maxwell 模型或空间离散误差足够小。固定背景分辨率、edge-space consistency、sub-cell geometry/source representation、材料模型和热空间 residual 仍需单独做收敛与物理验证。